% =====================================================================
%  chunks/rev_ch04.tex —— 第 4 章 读书笔记（tr-yi 撰）
% =====================================================================
\chapterreview{%
第 4 章是全书的目标章。前三章造的机器——同调、上同调、障碍、谱序列——到这里全部对准一个问题：光滑流形能否分类，能分到什么程度。章标题只写「光滑流形」，内容却远超名称：§1 至 §3 积攒不变量，§4 用手术把不变量变成分类定理，§5 又把高维机器降回三维、四维与 Hermitian $K$ 理论，让读者看见机器在哪里失效。

五节是五条主线。§1 给语言：联络、曲率、示性类与分类空间，示性类作为从向量丛范畴到上同调的自然变换由定义 1.1 给出，再与 $H^*(BG)$ 中的元素一一对应；(1.22) 说明曲率就是联络的换位子，即物理里的规范场强。§2 把第 3 章的同调整套搬到光滑范畴：de Rham 复形 (2.15) 与 de Rham 定理（该定理在书中无编号）、Hodge 分解 (2.24)，以及 (2.26) 借 Laplace 算子把解析挠率接回 Reidemeister 挠。§3 让流形自己充当同调类：Thom 同构 (3.9) 把 $H^j(B)$ 与 $H^{j+n}(T_\eta)$ 对齐，配边环 (3.8) 由此可算。§4 才谈分类：Smale--Hirsch 定理 (4.1)(4.2)、Kervaire--Milnor 的 $\Theta^n\to\pi_{N+n}(S^N)/J\pi_n(SO)$ (4.5)、刻画 PL 与光滑之差的 $\Gamma^n$（书中只给定义与性质，无编号公式），以及把同伦核换成同调核的手术 (4.8)--(4.10)。§5 回到低维与代数：Novikov 符号差公式 (5.3)、$s$-配边的 Whitehead 群障碍 (5.12)、Hermitian $K$ 理论。这个顺序不是随意的：先有不变量才谈分类，末节则说明高维这套办法到四维便断。

三条结论最该记住。其一，Thom 同构 (3.9) 把「用子流形表示同调类」化为可算的代数，配边论从此与同调论并列，而 (3.10)--(3.12) 又把 Stiefel--Whitney 类写成 Steenrod 运算的拉回。其二，符号差：$\tau(M^4)=p_1/3$ (3.13) 与可加性 (3.14) 表明它是配边不变量；到 §5 变成 $\langle L_k(p_1,\dots,p_k),z\rangle=\tau(M^{4k})$ (5.3)，正是 Novikov 1966 年证明有理 Pontryagin 类为拓扑不变量的入口。其三，$h$-配边定理使 $\Theta^n$ 在 $n\ge5$ 时精确分类同伦球面 (4.5)，怪球问题于是变成稳定同伦群的计算。

四处易失足。（1） $h$-配边定理给的是微分同胚，不是同伦等价；反过来同伦等价推不出同胚，透镜空间即是反例。（2） 配边强于同调：同调类可由子流形表示要 Thom 的判据，只有 $\Omega^{SO}_*\otimes\Q$ 才与多项式环同构，两者不可混同。（3） $s$-配边定理要求 $n\ge5$，且障碍落在 $Wh(\pi)=K_1(\pi)/\pm\pi$ (5.12)；四维的拓扑版本属 Freedman 的工作，本书未覆盖。（4） Kervaire 不变量（§4 讨论，无编号公式）只在 $n\equiv2\pmod4$ 可能非零、且取 $\Z/2$ 值，把它当作一切维数的障碍是常见误读；(3.5)(3.6) 是 Arf 恒等式与完全 Arf 不变量，二者不是一回事。

写得最扎实的是 §3，尤其 Thom 空间与 Milnor 怪球的构造；讲得最快的是 §2 的 Hodge 理论（(2.24) 后直接跳到 (2.26) 的解析挠率）与 §5 的 Hermitian $K$（只给定义与结果，Wall 群的谱序列计算全略）。Novikov 个人印记最重处也在 §5：环带猜想的 Kirby 路线、Mishchenko 的 $\tau^{\pi}$、以及把手术障碍译成 $\Z[\tfrac12]$ 上的 Hermitian $K$ 周期理论。若只带走一句话：本章的主题不是「光滑流形」，而是「分类的极限」——每多一层结构，不变量就多一层，而层与层之差正是 $\Theta^n$、$\Gamma^n$、$Wh(\pi)$ 这些有限群；读懂它们，才算读懂本章。
}

% CHUNK rev_ch04 DONE
